\section{\textsf{Functional Equations}}
\subsection{Problems}
\subsubsection{Kyrgzstan 2012 02/13/26}
\begin{problem}
Find all functions $ f:\RR\to\RR $ such that
\[
	f(f(x)^2+f(y)) = xf(x)+y \quad \forall x,y\in \RR
\]
\end{problem}
We claim that $f(x)=\pm x$ is the only possible solution. It is clear to see that $f(x)=\pm x$ indeed works.
Now we prove it is the only two functions that actually work.

\begin{claim}
	The function $f$ is bijective.
\end{claim}
\begin{proof}
	Plugging $x=0$ gives
	\[
		f(f(0)^2+f(0))=0,
	\]
	and plugging $u=f(0)^2+f(0)$ gives $f(u)=0$. Hence, we can say $x=u$ to find that
	\[
		f(f(y))=y. \qedhere
	\]
\end{proof}
Say $f(t)=x$. Hence,
\[
	f(t^2+f(y))=f(t)\cdot t + y= f(f(t)^2+f(y)) \iff f(t)^2=t^2.
\]
However, we must still escape the pointwise trap!

By the sake of contradiction, say the function $f$ isn't one of the two linear functions we found out. Hence, say
$f(a)=+a$ and $f(b)=-b$. Thus,
\[
	f(a^2-b)=a^2+b.
\]
Contradiction! $a$ and $b$ aren't equal 0 at the same time.


\clearpage

\subsubsection{02/14/26}
\begin{problem}
Find all functions $f:\mathbb{R}\to\mathbb{R}$ such that
\[
	f(x^{2}+y)=f(x^{27}+2y)+f(x^{4})
\]
for all $x,y\in\mathbb{R}$.
\end{problem}
We claim the answer is $f(x)=0$. Clearly it works, so we demonstrate now that it is the only possible answer.

Say $x^2+y=x^{27}+2y$. Hence,
\[
	f(2x^2-x^{27})=f(2x^2-x^{27})+f(x^4) \iff f(x^4)=0.
\]
Thus, $f(x)=0$ for all nonnegative numbers. To prove the same is true for negative numbers and $0$,
we must simply say $y=0$, resulting in $f(x^{27})=0$.

\clearpage

\subsubsection{02/14/26}
\begin{problem}
Find all strictly increasing functions $f:\mathbb{Z}\to\mathbb{Z}$ such that
\[
	f(f(x))=x+2
\]
for all integers $x$.
\end{problem}
